% 奥利弗·亥维赛（Oliver Heaviside）（综述）
% license CCBYSA3
% type Wiki

本文根据 CC-BY-SA 协议转载翻译自维基百科\href{https://en.wikipedia.org/wiki/Oliver_Heaviside}{相关文章}

\begin{figure}[ht]
\centering
\includegraphics[width=6cm]{./figures/1b91699cfa36da0e.png}
\caption{约1900年的赫维赛德} \label{fig_ALFkws_3}
\end{figure}
奥利弗·赫维赛德（Oliver Heaviside，/ˈhɛvisaɪd/，读作“赫维赛德”\(^\text{[2]}\)，1850年5月18日－1925年2月3日）是一位英国数学家和电气工程师。他发明了一种求解微分方程的新方法（与拉普拉斯变换等价），独立发展了向量分析，并将麦克斯韦方程组改写成如今通用的形式。他对麦克斯韦方程组在麦克斯韦逝世后的理解与应用产生了深远影响。此外，1893年，赫维赛德将麦克斯韦方程扩展到了“引力电磁学”领域——这一理论在2005年的“重力探测器B”实验中得到了验证。他提出的“电报方程”在他生前的晚期获得了商业上的重要意义，尽管这一成果曾长期未被重视，因为当时很少有人理解他所采用的创新数学方法。\(^\text{[3]}\)尽管赫维赛德一生大部分时间都与主流科学界格格不入，但他彻底改变了电信、数学和科学的面貌。\(^\text{[3]}\)
\subsection{早年经历}
奥利弗·赫维赛德于1850年5月18日出生在英国伦敦卡姆登镇国王街55号（现为普伦德街，Plender Street）\(^\text{[4]: 13 }\)，是托马斯·赫维赛德和雷切尔·伊丽莎白·韦斯特的第三个孩子，也是家中最小的儿子。他的父亲是一名制图员和木版雕刻师。赫维赛德从小个子矮小、头发红亮，年幼时患上猩红热，导致听力受损。家里继承了一小笔遗产，使得全家在他十三岁时搬到了卡姆登较好的地区，他也因此得以进入卡姆登文法学校就读。赫维赛德是个出色的学生，1865年在500名学生中名列第五。然而，由于家庭经济拮据，他在16岁时被迫辍学。之后，他自学了一年，并未再接受任何正式教育。\(^\text{[5]: 51 }\)

赫维赛德的姨夫是查尔斯·惠斯通爵士，这位国际知名的电报与电磁学专家是19世纪30年代中期首个商业化电报系统的共同发明者。惠斯通非常关心外甥的教育，\(^\text{[6]}\)并于1867年将赫维赛德送往英国纽卡斯尔（，在那里，他在惠斯通的哥哥亚瑟·惠斯通管理的一家电报公司任职。\(^\text{[5]: 53 }\)

两年后，赫维赛德在丹麦“大北电报公司”找到了一份电报操作员的工作，负责从纽卡斯尔铺设通往丹麦的海底电缆，施工方为英国承包商。不久之后，他被提升为电气工程师。在工作期间，赫维赛德仍坚持自学。22岁时，他在著名的《哲学杂志》上发表了论文《惠斯通电桥在给定电流计与电池条件下测量特定电阻的最佳配置》\(^\text{[7]}\)。这篇论文得到了多位物理学家的高度评价——其中包括威廉·汤姆孙爵士（Sir William Thomson，后来的开尔文勋爵）和詹姆斯·克拉克·麦克斯韦，他们此前都曾尝试解决这一代数问题但未成功。赫维赛德还亲自将论文送给了汤姆孙爵士。此后，他又发表了一篇关于“双工传输”（在电报电缆中的应用的论文\(^\text{[8]}\)，并在文中讽刺了当时英国邮政电报系统总工程师R.S.卡利，后者一直否认双工电报的可行性。1873年晚些时候，赫维赛德申请加入“电报工程师学会”，却被拒绝，理由是“他们不希望接收电报职员”。这一侮辱激怒了赫维赛德，他于是请威廉·汤姆孙为其作推荐人，并在该学会会长的支持下，最终“尽管邮局里那些势利小人反对”，他还是成功入会。\(^\text{[5]: 60 }\)

1873年，赫维赛德读到了刚刚出版、后来享誉世界的詹姆斯·克拉克·麦克斯韦两卷本巨著《电与磁论》。多年以后，他在晚年回忆道：

“我还记得年轻时第一次看到麦克斯韦那部伟大的著作……我立刻意识到它是伟大的、更加伟大的、最伟大的，其力量中蕴藏着惊人的潜能……我下定决心要精通这本书，并立即开始钻研。当时我非常无知，对数学分析一无所知（只学过一些中学代数和三角学，还大多忘了），因此我知道自己该怎么补课。我花了好几年时间，才尽我所能去理解其中的内容。之后，我暂且放下麦克斯韦，按照自己的方式继续研究。那时我的进步要快得多……可以说，我所宣扬的，是我自己对麦克斯韦理论的理解与诠释。”\(^\text{[8]}\)

赫维赛德在家中独自从事研究，发展出了传输线理论（transmission line theory，也称为“电报方程”）。他通过数学推导证明：在电报线路中，如果沿线均匀分布电感，就可以减小信号的衰减和失真；并且，如果电感足够大且绝缘电阻不过高，那么电路将不会发生频率失真——也就是说，各频率的电流将以相同的速度传播。\(^\text{[10]}\)赫维赛德的方程为电报系统的改进和推广提供了重要的理论基础。
\subsection{中年时期}
从1882年至1902年（中间有三年例外），赫维赛德定期为行业刊物《电气工程师》撰写文章。该杂志希望借助他的学术声望提升自身影响力，因此每年支付他40英镑稿酬。虽然这笔钱几乎不足以维生，但赫维赛德生活要求极低，而且这份工作正是他最热爱的事情。1883年至1887年间，他平均每月发表2至3篇论文，这些文章后来构成了他的重要著作《电磁理论》和《电学论文集》的主要内容。\(^\text{[5]: 71 }\)

1880年，赫维赛德研究了电报传输线中的“集肤效应”，同年他在英国申请了“同轴电缆”的专利。1884年，他将麦克斯韦原本复杂的数学表达式（当时曾以四元数形式改写）重新整理为现代的“向量表示法”，将原先20个未知数的20个方程大大简化为如今广为人知的4个偏微分方程，也就是我们今天所说的“麦克斯韦方程组”。这四个方程系统地描述了电荷（静态与动态）、磁场及其相互关系——即电磁场的本质。

1880年至1887年间，赫维赛德发展了“算子演算”，以符号$p$（布尔此前使用$D$表示\(^\text{[11]}\)）代表微分算子，从而可以将微分方程直接转化为代数方程求解。这种方法后来因“缺乏数学严谨性”而引发广泛争议。赫维赛德对此曾著名地辩解道：“数学是一门实验科学，定义并不是先行的，而是在研究对象的本质逐渐显现后才自然形成的。”\(^\text{[12]}\)他还幽默地反问过一句：“难道仅仅因为我不能完全理解消化的过程，就要拒绝吃饭吗？”\(^\text{[13]}\)1887年，赫维赛德与他的哥哥亚瑟合作撰写了一篇题为《桥式电话系统》的论文。然而，这篇论文被亚瑟的上司——英国邮政局的威廉·亨利·普里斯阻止发表。原因是论文中提出应在电话和电报线路中加入“负载线圈”（loading coils，即电感器），以增强自感，从而校正信号传输中的失真。而普里斯此前刚刚宣称“自感是清晰传输的最大敌人”，因此他坚决反对。\(^\text{[14]}\)赫维赛德还确信，正是普里斯在背后操纵，迫使《电气工程师》的编辑被解雇，导致他连载多年的论文系列被迫中断（直到1891年才恢复）。赫维赛德与普里斯之间的敌意由来已久。赫维赛德认为普里斯在数学上“完全不懂行”，传记作者保罗·J·纳欣也持相同观点：“普里斯是一个掌握权力的政府官员，野心勃勃，但在某些令人惊讶的方面却是个十足的笨蛋。” 普里斯阻挠赫维赛德工作的动机，更多是为了维护自己的声誉、避免承认错误，而非出于对赫维赛德研究本身的质疑。\(^\text{[4]: xi–xvii, 162–183 }\)

赫维赛德的研究成果在发表后很长一段时间内都未被重视。直到1897年，美国电话电报公司才开始关注这项工作，并聘请公司科学家乔治·A·坎贝尔以及外部研究员迈克尔·I·普平对赫维赛德的理论进行验证，试图找出其中的不足或错误。坎贝尔和普平在赫维赛德的基础上进行了扩展研究，AT&T随即申请了多项专利，不仅包括他们自己的研究成果，还涵盖了赫维赛德此前发明的线圈构造方法。后来，AT&T向赫维赛德提出金钱补偿，以换取其相关权利。据说这一提议部分源于贝尔系统工程师们对赫维赛德的尊敬。然而，赫维赛德拒绝了这一提议，表示除非公司公开承认其发明，否则他不会接受任何报酬。尽管他长期生活贫困，但他仍坚定地拒绝了这笔钱——这一举动尤为令人钦佩。1959年，诺伯特·维纳发表了小说《诱惑者》，在书中影射性地指控AT&T（化名为“威廉斯控制公司”Williams Controls Company）和迈克尔·I·普平（化名为“迭戈·多明格斯”Diego Dominguez）窃取了赫维赛德的发明成果。\(^\text{[15][16][17]}\)

然而，这次挫折使赫维赛德将注意力转向了电磁辐射。\(^\text{[18]}\)在他于1888年和1889年发表的两篇论文中，他计算了运动电荷周围电场与磁场的变形，以及电荷进入更高密度介质时所产生的效应。这项研究预言了后来被称为契伦科夫辐射的现象，并启发了他的好友乔治·菲茨杰拉德提出了如今所知的洛伦兹–菲茨杰拉德收缩假设。

1889年，赫维赛德首次发表了运动带电粒子所受磁力的正确推导，\(^\text{[19]}\)这正是现在所称洛伦兹力的磁场分量部分。

在19世纪80年代末至90年代初，赫维赛德进一步研究了电磁质量的概念。他认为电磁质量可被视为一种“物质质量”，能够产生相同的物理效应。后来，威廉·维恩在低速条件下验证了赫维赛德的公式。

1891年，英国皇家学会以授予他院士称号的方式，正式表彰他在电磁现象数学化描述方面的杰出贡献。次年，《皇家学会哲学汇刊》专门刊登了超过五十页篇幅的论文，介绍他的向量分析方法与电磁理论。

1894年，他被选为曼彻斯特文学与哲学学会的名誉会员。\(^\text{[20]}\)1905年，德国哥廷根大学授予他荣誉博士学位，以表彰其在电磁学方面的卓越成就。
\subsection{晚年与思想观点}
\begin{figure}[ht]
\centering
\includegraphics[width=6cm]{./figures/df06c8c7ba5b272d.png}
\caption{佩恩顿为纪念赫维赛德而设立的蓝色纪念牌} \label{fig_ALFkws_1}
\end{figure}
1896年，乔治·菲茨杰拉德和约翰·佩里为赫维赛德争取到了一项政府养老金，每年金额为120英镑。当时他已定居在德文郡。他们劝说他接受这笔资助，因为赫维赛德此前曾拒绝过皇家学会提供的其他慈善援助。\(^\text{[18]}\)

1902年，赫维赛德提出了如今被称为肯内利–赫维赛德层的电离层概念。他的理论解释了无线电信号能够绕地球曲率传播的原理。电离层的存在后来于1923年得到证实。赫维赛德的预言，加上普朗克辐射理论，可能使人们在此后的几十年中，对探测来自太阳及其他天体的无线电波失去了兴趣。无论出于何种原因，之后的三十年间几乎无人尝试，直到1932年詹斯基开创了射电天文学。

赫维赛德是阿尔伯特·爱因斯坦相对论的公开反对者。\(^\text{[21]}\)数学家霍华德·伊夫斯评论说，赫维赛德“是当时唯一一个质疑爱因斯坦的顶尖物理学家，他对相对论的批评常常近乎荒谬”。\(^\text{[21]}\)

在晚年，赫维赛德的行为变得相当古怪。据同事B.A.贝伦德回忆，他逐渐与世隔绝，厌恶与人接触——甚至把自己为《电气工程师》撰写的论文稿件送到杂货店，由编辑亲自前去取走。\(^\text{[22]}\)年轻时的赫维赛德是个热衷骑行的人，但到了六十岁后，他的健康状况严重恶化。在此期间，他常在签名后附上缩写“W.O.R.M.”（意为“虫”）。据传，他还开始将指甲涂成粉红色，并搬来几块花岗岩石块当作家具使用。\(^\text{[4]: xx }\)1922年，他成为首届法拉第奖章的获得者，该奖正是在那一年设立。

在宗教观方面，赫维赛德自称是一神论者，但并不虔诚。据说他还曾嘲笑那些相信至高神明的人。\(^\text{[23]}\)

1925年2月3日，赫维赛德在英国托基因从梯子上跌落去世，享年74岁。\(^\text{[24]}\)他的墓地位于佩因顿公墓东南角附近建筑物的后方偏右处，与父亲托马斯和母亲雷切尔合葬。墓碑于2005年由一位匿名捐助者出资清理修复。\(^\text{[25]}\)赫维赛德在电气工程师群体中始终备受尊敬，尤其是在他对开尔文传输线分析的修正得到证实之后；不过，他在更广泛公众中的声誉主要是在身后才真正建立起来的。
\subsubsection{赫维赛德纪念计划}
\begin{figure}[ht]
\centering
\includegraphics[width=6cm]{./figures/5b52f8d75ceb08ea.png}
\caption{修复项目实施前后的对比} \label{fig_ALFkws_2}
\end{figure}
2014年7月，英国纽卡斯尔大学与“纽卡斯尔电磁学兴趣小组”的学者共同发起了“赫维赛德纪念计划”，旨在通过公众募捐的方式对赫维赛德纪念碑进行全面修复。\(^\text{[26][27][28]}\)修复后的纪念碑于2014年8月30日举行揭幕仪式，由赫维赛德的远亲阿兰·希瑟主持揭幕。出席仪式的嘉宾包括托贝市市长、托贝地区国会议员、代表英国工程与技术学会的前科学博物馆馆员、托贝市民协会主席，以及来自纽卡斯尔大学的代表。\(^\text{[29]}\)
\subsubsection{英国工程与技术学会}
赫维赛德的部分论文与资料现由英国工程与技术学会档案中心收藏。\(^\text{[30]}\)该收藏包括：装满数学方程与计算的笔记本、带有批注的小册子（主要与电报技术相关）、手稿笔记、论文草稿、通信往来，以及《电磁理论》的文章初稿。此外，还保存了一份1950年由贝尔电话实验室总裁奥利弗·E·巴克利录制的赫维赛德致敬音频，该音频已被数字化并可通过IET档案馆的“奥利弗·赫维赛德传记”页面在线访问。\(^\text{[31]}\)

1908年，赫维赛德被授予英国电气工程师学会荣誉会员称号。他在《IEE荣誉会员与法拉第奖章获得者名录（1871–1921）》中的登记为：“1908 Oliver Heaviside FRS”。[32][33]1922年，他成为新设立的法拉第奖章的首位获奖者。1950年，英国电气工程师学会理事会设立了赫维赛德优等奖：“委员会已考虑为奥利弗·赫维赛德建立一种永久性纪念形式，因此建议设立‘赫维赛德优等奖’，每年奖励10英镑，授予当年录用的最佳数学论文作者。”\(^\text{[34]}\)
\subsection{创新与发现}
赫维赛德在向量方法与向量分析的发展与推广方面作出了重大贡献。\(^\text{[35]}\)麦克斯韦最初的电磁学表述包含20个方程与20个变量。赫维赛德引入了向量分析中的旋度与散度算子，将其中 12 个方程重新推导为4个方程、4个变量（$\mathbf{B}, \mathbf{E}, \mathbf{J}, \rho$），也就是我们今天熟知的麦克斯韦方程组形式。值得注意的是，赫维赛德的方程与麦克斯韦原始形式并不完全相同，事实上，赫维赛德的表述更容易被修改以适应量子物理的框架。\(^\text{[36]}\)赫维赛德还利用引力与电学中平方反比定律的类比，讨论了引力波存在的可能性。\(^\text{[37]}\)在四元数乘法中，向量的平方为负值，这令赫维赛德十分不满。他主张取消这种负号，因此被C.J.乔利认为是双曲四元数的创始人之一，尽管这种数学结构的主要贡献应归功于亚历山大·麦克法兰。\(^\text{[38]}\)

赫维赛德发明了著名的赫维赛德阶跃函数，用于计算电路接通时的电流变化。他也是第一个使用单位脉冲函数（即如今所称的狄拉克$\delta$函数，Dirac delta function）的人。\(^\text{[39]}\)此外，他发明了用于求解线性微分方程的算子演算方法。这一方法与后来基于布罗姆维奇积分的拉普拉斯变换方法极为相似；后者由布罗姆维奇提出，用复平面路径积分的方式对赫维赛德的方法进行了严格的数学化证明。\(^\text{[40]}\)赫维赛德本人熟悉拉普拉斯变换，但他认为自己的方法更加直接、简洁。\(^\text{[41][42]}\)

赫维赛德提出了传输线理论（transmission line theory，又称电报方程telegrapher’s equations），这一理论使跨大西洋电缆的传输速率提高了十倍。在此之前，每传送一个字符需要十分钟，而应用该理论后，速度立即提升到每分钟一个字符。与此密切相关的是，他还发现：在电话线路中串联电感可以显著改善传输质量。\(^\text{[43]}\)赫维赛德还独立发现了坡印廷矢量的概念，用以描述电磁能量在空间中的流动。\(^\text{[4]: 116–118 }\)

此外，赫维赛德提出：地球高层大气中存在一个电离层，即后来的电离层。他预言了后来被称为肯内利–赫维赛德层的存在。1947年，爱德华·阿普尔顿因通过实验证实这一电离层的存在，而获得了诺贝尔物理学奖。
\subsubsection{电磁学术语}
赫维赛德在电磁理论中创造了以下一系列专业术语：
\begin{itemize}
\item Admittance（导纳）：阻抗的倒数（1887年12月提出）；
\item Elastance（电弹性）：电容的倒数，即容许度的倒数（1886年提出）；
\item Conductance（电导）：导纳的实部，即电阻的倒数（1885年9月提出）；
\item Electret（驻极体）：电学上对应于永久磁体的概念，即具有准永久电极化（如铁电材料）的物质；
\item Impedance（阻抗）（1886年7月提出）；
\item Inductance（电感）（1886年2月提出）；
\item Permeability（磁导率）（1885年9月提出）；
\item Permittance（容许度）与 Permittivity（介电常数）（1887年6月提出）；
\item Reluctance（磁阻）（1888年5月提出）。[44]
\end{itemize}
此外，赫维赛德常被错误地认为是以下两个术语的创造者：Susceptance（电纳）——导纳的虚部；实际上由查尔斯·普罗透斯·斯坦梅茨于1894年提出。\(^\text{[45]}\)Reactance（电抗）——阻抗的虚部；实际上由爱德华·奥斯皮塔利耶于1893年提出。\(^\text{[46]}\)
\subsection{出版著作}
\begin{itemize}
\item 1885、1886、1887年 — Electromagnetic Induction and its Propagation（《电磁感应及其传播》），发表于《The Electrician》。
\item 1888–1889年 — Electromagnetic Waves, the Propagation of Potential, and the Electromagnetic Effects of a Moving Charge（《电磁波、电势传播与运动电荷的电磁效应》），发表于《The Electrician》。
\item 1889年 — On the Electromagnetic Effects due to the Motion of Electrification through a Dielectric*（《电介质中电荷运动引起的电磁效应》），发表于《Philosophical Magazine》，第5辑第27卷，第324页。
\item 1892年 — On the Forces, Stresses, and Fluxes of Energy in the Electromagnetic Field（《电磁场中的力、应力与能量流》），发表于《Philosophical Transactions of the Royal Society A》，第183卷，页423–480。
\item 1892年 — On Operators in Physical Mathematics, Part I（《物理数学中的算子·第一部分》），发表于《Proceedings of the Royal Society》，1892年1月1日，第52卷，页504–529。
\item 1892年 — 《Electrical Papers》（《电学论文集》），第1卷，伦敦与纽约：麦克米伦公司。ISBN 9780828402354。
\item 1893年 — On Operators in Physical Mathematics, Part II（《物理数学中的算子·第二部分》），发表于《Proceedings of the Royal Society》，1893年1月1日，第54卷，页105–143。
\item 1893年 — A Gravitational and Electromagnetic Analogy（《引力与电磁类比》），发表于《The Electrician》，第31卷，第281–282页（第一部分），第359页（第二部分）。
\item 收录于《Electromagnetic Theory》（《电磁理论》）第1卷第4章附录B，第455–466页。
\item 1893年 — 《Electromagnetic Theory》（《电磁理论》），第1卷，伦敦：The Electrician Printing and Publishing Co.。ISBN 978-0-8284-0235-4。[47]
\item 1894年 — 《Electrical Papers》（《电学论文集》），第2卷，伦敦与纽约：麦克米伦公司（Macmillan Co.）。
\item 1899年 — 《Electromagnetic Theory》（《电磁理论》），第2卷，伦敦：The Electrician Printing and Publishing Co.。
\item 1912年 — 《Electromagnetic Theory》（《电磁理论》），第3卷，伦敦：The Electrician Printing and Publishing Co.。
\item 1925年 — 《Electrical Papers》（《电学论文集》），共2卷，波士顿：Copley出版社。
\item 1950年 — 《Electromagnetic Theory: The Complete & Unabridged Edition》（《电磁理论：完整未删节版》），Spon出版社出版，1950年由Dover再版。
\item 1970年 — 《Electrical Papers》（《电学论文集》），纽约：切尔西出版社。ISBN 978-0-8284-0235-4。
\item 1971年 — 《Electromagnetic Theory: Including an Account of Heaviside's Unpublished Notes for a Fourth Volume》（《电磁理论：含赫维赛德第四卷未发表笔记》），切尔西出版社，ISBN 0-8284-0237-X。
\item 2001年 — 《Electrical Papers》（《电学论文集》），美国数学学会出版。ISBN 978-0-8218-2840-3。
\end{itemize}
\subsection{另见}
\begin{itemize}
\item 1850年的科学发展
\item 电位移场
\item 毕奥–萨伐尔定律
\item 电桥电路 § 赫维赛德电桥
\end{itemize}
\subsection{参考文献}
\begin{enumerate}
\item Anon（1926）. “讣告：已故院士名单——鲁道夫·梅塞尔、弗雷德里克·托马斯·特劳顿、约翰·范恩、约翰·杨·布坎南、奥利弗·赫维赛德、安德鲁·格雷。” 《皇家学会学报A：数学、物理与工程科学》，第110卷（756期），页 i–v。Bibcode:1926RSPSA.110D...1.. doi:10.1098/rspa.1926.0036。
\item “Heaviside.” dictionary.com。Dictionary.com。检索于2025年10月8日。
\item Hunt, B. J.（2012）. “奥利弗·赫维赛德：一个杰出的怪才”。Physics Today，第65卷第11期，页48–54。Bibcode:2012PhT....65k..48H。doi:10.1063/PT.3.1788。
\item Nahin, Paul J.（2002年10月9日）. Oliver Heaviside: The Life, Work, and Times of an Electrical Genius of the Victorian Age（《奥利弗·赫维赛德：维多利亚时代电气天才的一生》）。约翰·霍普金斯大学出版社（JHU Press）。ISBN 978-0-8018-6909-9。
\item Bruce J. Hunt（1991）. The Maxwellians（《麦克斯韦学派》）。康奈尔大学出版社（Cornell University Press）。ISBN 978-0-8014-8234-2。
\item Sarkar, T. K.; Mailloux, Robert; Oliner, Arthur A.; Salazar-Palma, M.; Sengupta, Dipak L.（2006）. History of Wireless（《无线通信史》）。约翰·威立父子公司（John Wiley & Sons）。第230页。ISBN 978-0-471-78301-5。
\item Heaviside（1892），页3–8。
\item Heaviside（1892），页18–34。
\item Sarkar, T. K. 等（2006年1月30日）. History of Wireless。约翰·威立父子公司，第232页。ISBN 978-0-471-78301-5。
\item 以下内容部分引自现已进入公共领域的出版物：Kempe, Harry Robert（1911），“Telephone（电话）”。载于 Hugh Chisholm（编），Encyclopædia Britannica，第26卷（第11版），剑桥大学出版社，第554页。
\item A Treatise on Differential Equations（《微分方程论》），1859年。
\item “VIII. On Operations in Physical Mathematics. Part II.” Proceedings of the Royal Society of London，第54卷（326–330期），页105–143，1894年。doi:10.1098/rspl.1893.0059。S2CID 121790063。
\item Heaviside，《Mathematics and the Age of the Earth》（《数学与地球的年龄》），收录于《电磁理论》第2卷。
\item Hunt, Bruce J.（2004）. “Heaviside, Oliver.” Oxford Dictionary of National Biography（《牛津国家人物传记词典》），在线版，牛津大学出版社。doi:10.1093/ref:odnb/33796。（需订阅，或通过Wikipedia Library或英国公共图书馆访问。）
\item Wiener, Norbert（1993）. Invention: The Care and Feeding of Ideas（《发明：思想的养成与滋养》）。美国马萨诸塞州剑桥：麻省理工学院出版社（MIT Press）。页70–75。ISBN 0-262-73111-8。
\item Wiener, Norbert（1959）. The Tempter（《诱惑者》）。纽约：兰登书屋（Random House）。
\item Montagnini, Leone（2017）. Harmonies of Disorder – Norbert Wiener: A Mathematician-Philosopher of Our Time（《混沌的和谐——诺伯特·维纳：我们时代的数学哲学家》）。瑞士楚格：施普林格出版社（Springer）。页249–252。ISBN 978-3-31984455-8。
\item Hunt（2004）。
\item Heaviside, O.（1889）. “XXXIX. On the Electromagnetic Effects due to the Motion of Electrification through a Dielectric”（《电介质中电荷运动引起的电磁效应》）。Philosophical Magazine，第五辑，第27卷（167期），页324–339。doi:10.1080/14786448908628362。
\item Memoirs and Proceedings of the Manchester Literary & Philosophical Society（《曼彻斯特文学与哲学学会论文集》），第四辑，第八卷，1894年。
\item Eves, Howard（1988）. 《重返数学圈：第五辑数学故事与轶事集》。PWS-Kent 出版公司，第27页。ISBN 9780871501059。
\item “Pages with the Editor”（编辑寄语）（PDF），刊于《Popular Radio》（《大众无线电》），第7卷第6期，纽约，1925年6月，第6页。检索于2014年8月14日。
\item Pickover, Clifford A.（1998）. 《奥利弗·赫维赛德》，载于《奇异的头脑与天才：古怪科学家与疯子的秘密生活》。Plenum 出版公司。ISBN 9780306457845。
宗教信仰：一神论者，但并不虔诚；曾嘲笑那些信奉至高存在的人。
\item “Oliver Heaviside”，Journal of the AIEE（《美国电气工程师学会期刊》）讣告，卷44，第3期，页316–317，1925年3月。doi:10.1109/JAIEE.1925.6537168。S2CID 51663331。
\item Mahon, Basil（2009）. Oliver Heaviside: Maverick Mastermind of Electricity（《奥利弗·赫维赛德：电学的特立独行天才》）。英国工程与技术学会（The Institution of Engineering and Technology）。ISBN 9780863419652。
\item “Heaviside Memorial Project Homepage.” Nature，第165卷（4208期），《赫维赛德纪念计划》，页991–993，2014年7月27日。原文存档于2014年7月18日，检索于2014年7月31日。
\item “Bid to Restore Paignton Monument to Oliver Heaviside.” *Herald Express*，网站：[www.torquayheraldexpress.co.uk，2014年7月27日。原文存档于2014年8月6日，检索于2014年7月29日。](http://www.torquayheraldexpress.co.uk，2014年7月27日。原文存档于2014年8月6日，检索于2014年7月29日。)
\item “The Heaviside Memorial Project.” Newcastle University，网站：[www.newcastle.ac.uk，2014年7月29日。原文存档于2014年7月29日，检索于同日。](http://www.newcastle.ac.uk，2014年7月29日。原文存档于2014年7月29日，检索于同日。)
\item “Restored Heaviside Memorial Unveiled on Saturday.” Herald Express，网站：[www.torquayheraldexpress.co.uk，2014年9月1日。原文存档于2014年9月3日，检索于2014年9月1日。](http://www.torquayheraldexpress.co.uk，2014年9月1日。原文存档于2014年9月3日，检索于2014年9月1日。)
\item Savoy Hill House，7–10 Savoy Hill，London WC2R 0BU。电子邮件：[archives@theiet.org](mailto:archives@theiet.org)。
\item Heaviside, Oliver. “Oliver Heaviside 1850–1925.” The IET Archives: Biographies（IET 档案馆传记）。检索于2023年11月21日。
\item Heaviside, Oliver（1908）. “IEE Roll of Honorary Members and Faraday Medallists 1871–1921.” IET Archives Reference:*IET/SPE/4/8/1。
\item Heaviside, Oliver（2023年4月5日）. “From under the sea to the edge of space: The work of Oliver Heaviside.” Stories of the Institution of Electrical Engineers. The IET Archives Blog: Stories from the Institution of Engineering and Technology. 检索于2023年11月21日。
\item Heaviside Premium Award（1950年2月2日）. “IEE Council Minutes.” IET Archive Reference: IET/ORG/2/1/19。
\item 参见特别是《电磁理论》（Electromagnetic Theory，1893）中“向量代数与分析的要素”，第1卷第3章，第132–305页，赫维赛德在此完整阐述了现代体系。
\item Topological Foundations of Electromagnetism（《电磁学的拓扑基础》），《世界科学当代化学物理学系列》，2008年3月13日，作者：Terence W. Barrett。
\item 《电磁理论》（Electromagnetic Theory，1893），附录B《引力与电磁类比》，第455–466页。此文发表于爱因斯坦关于同类主题论文的25年前。
\item Hamilton（1899），由 C.J. Joly 编辑，《Elements of Quaternions》（《四元数原理》，第二版），Longmans, Green & Co.出版社，第163页。ISBN 9780828402194。
\item 《电磁理论》第2卷，第271段，方程54与55。
\item 参见Bromwich W.P.文献中引用的Jeffreys论文。
\item 《电磁理论》第3卷，自第324页起的章节，可在线阅读。
\item 赫维赛德算子演算的严格数学化版本见：Mikusinski, J.（1959），The Operational Calculus（《算子演算》），Pergamon Press。
\item Wiener, Norbert（1993）. Invention: The Care and Feeding of Ideas（《发明：思想的养成与滋养》），麻省理工学院出版社，页70–75。ISBN 0-262-73111-8。
\item Ronald R. Kline，《Steinmetz: Engineer and Socialist》（《斯坦梅茨：工程师与社会主义者》），第337页，约翰·霍普金斯大学出版社，1992年。ISBN 0801842980。
\item Kline，第88页。
\item Steinmetz, Charles Proteus；Bedell, Frederick（1894），“Reactance”，Transactions of the American Institute of Electrical Engineers（《美国电气工程师学会汇刊》），第11卷，第768–776页。
参见：Blondel, A.，“A propos de la reactance”（《论电抗》），刊于 L’Industrie Electrique，1893年5月10日。
赫维赛德本人也确认：“术语‘reactance’最近在法国被提出，我认为它是一个非常实用的词。”——见《电磁理论》，第1卷，第439页，1893年。
\item Swinburne, J.（1894）. “Review of Electromagnetic Theory, Vol. I”（《〈电磁理论〉第一卷书评》）。Nature，第51卷（1312期），页171–173。doi:10.1038/051171a0。S2CID 3940841。
\end{enumerate}
\subsection{延伸阅读}
\begin{itemize}
\item Philosophical Transactions of the Royal Society A: Mathematical, Physical and Engineering Sciences，《纪念奥利弗·赫维赛德〈电磁理论〉发表125周年》（主题专刊），第37卷，第2134期，2018年12月13日。
\item The Heaviside Centenary Volume（《赫维赛德百年纪念文集》），英国电气工程师学会，伦敦，1950年。
\item Berg, E. J.（1929）. Heaviside’s Operational Calculus as Applied to Engineering and Physics（《赫维赛德算子演算在工程与物理中的应用》）。麦格劳–希尔出版社（McGraw-Hill），《电气工程教材》系列。
\item Buchwald, Jed Z.（1985）. From Maxwell to Microphysics: Aspects of Electromagnetic Theory in the Last Quarter of the Nineteenth Century（《从麦克斯韦到微观物理：十九世纪后期电磁理论的若干方面》）。芝加哥大学出版社。ISBN 978-0-226-07882-3。
\item Calvert, James B.（2002）. *Heaviside, Laplace, and the Inversion Integral（《赫维赛德、拉普拉斯与反演积分》），丹佛大学。
\item Hunt, Bruce J.（1991）. *The Maxwellians（《麦克斯韦学派》），平装本版（2005年）。康奈尔大学出版社（Cornell University Press）。ISBN 978-0-8014-8234-2。
\item Jackson, W.（1950）. “Life and Work of Oliver Heaviside (May 18, 1850 – February 3, 1925)”（《奥利弗·赫维赛德的生平与工作》）。Nature，第165卷（4208期），1950年6月24日出版，页991–993。Bibcode:1950Natur.165..991J。doi:10.1038/165991a0。PMID 15439051。
\item effreys, Harold（1927）. Operational Methods in Mathematical Physics*（《数学物理中的算子方法》）。剑桥大学出版社，第二版，1931年。
\item Josephs, H. J.（1963）. Oliver Heaviside: A Biography（《奥利弗·赫维赛德传》）。伦敦出版。
\item Laithwaite, E. R.（1982）. “Oliver Heaviside – Establishment Shaker”（《奥利弗·赫维赛德——体制的撼动者》），Electrical Review，1982年11月12日。
\item Lee, G.（1947）. Oliver Heaviside（《奥利弗·赫维赛德》）。伦敦出版。
\item Lützen, J.（1980）. “Heaviside’s Operational Calculus and the Attempts to Rigorize It”（《赫维赛德算子演算及其数学化的尝试》），Archive for History of Exact Sciences，第21卷，页161–200。
\item Lynch, A. C.（1991）. 载于 G. Hollister-Short（编），“The Sources for a Biography of Oliver Heaviside”（《奥利弗·赫维赛德传记的资料来源》），History of Technology*（《技术史》），伦敦与纽约，第13卷。
\item Mahon, Basil（2009年5月11日）. Oliver Heaviside: Maverick Mastermind of Electricity（《奥利弗·赫维赛德：电学的特立独行天才》）。英国工程与技术学会（Institution of Engineering and Technology）。ISBN 978-0-86341-965-2。
\item Mende, F. F.（2015）. “What is Not Taken into Account and They Did Not Notice Ampere, Faraday, Maxwell, Heaviside and Hertz”（《被忽略的事实：安培、法拉第、麦克斯韦、赫维赛德与赫兹未被察觉的贡献》），AASCIT Journal of Physics，第1卷第1期，页28–52，2015年3月。
\item Moore, Douglas H.; Whittaker, Edmund Taylor（1928）. Heaviside Operational Calculus: An Elementary Foundation（《赫维赛德算子演算基础》）。美国爱思唯尔出版社（American Elsevier Publishing Company）。ISBN 0-444-00090-9。
\item Nahin, Paul J.（1987）. Oliver Heaviside, Sage in Solitude: The Life, Work, and Times of an Electrical Genius of the Victorian Age（《奥利弗·赫维赛德：孤独的贤者——维多利亚时代电气天才的一生》）。IEEE 出版。ISBN 978-0-87942-238-7。
\item Rocci, Alessio（2020）. “Back to the Roots of Vector and Tensor Calculus: Heaviside versus Gibbs”（《回溯向量与张量分析的起源：赫维赛德与吉布斯之争》），*Archive for History of Exact Sciences。doi:10.1007/s00407-020-00264-x。
\item Whittaker, E. T.（1929）. “Oliver Heaviside.” Bulletin of the Calcutta Mathematical Society，第20卷（1928–1929），页199–220。
\item Yavetz, I.（1995）. From Obscurity to Enigma: The Work of Oliver Heaviside, 1872–1889（《从默默无闻到科学谜团：赫维赛德的研究（1872–1889）》）。Birkhäuser出版社。ISBN 978-3-7643-5180-9。
\end{itemize}
\subsection{外部链接}
\begin{itemize}
\item 维基共享资源上的奥利弗·赫维赛德相关媒体
\item 迪布纳图书馆肖像收藏：“Oliver Heaviside”
\item 互联网档案馆上赫维赛德的著作与相关资料
\item Fleming, John Ambrose（1911）.“Units, Physical”（《物理单位》），载于《大英百科全书》第27卷（第11版），页738–745。
\item Ghigo, F. “Pre-History of Radio Astronomy, Oliver Heaviside (1850–1925)”（《射电天文学的前史：奥利弗·赫维赛德（1850–1925）》），美国国家射电天文台（NRAO），西弗吉尼亚州格林班克。原文存档于2020年6月15日。
\item Gustafson, Grant. “Heaviside’s Methods”（《赫维赛德的方法》），犹他大学数学系网站 math.Utah.edu。（PDF）
\item Heather, Alan. “Oliver Heaviside”，托贝业余无线电协会。原文存档于2017年9月24日。
\item Katz, Eugenii. “Oliver Heaviside”，耶路撒冷希伯来大学，原文存档于2009年10月27日。
\item Leinhard, John H.（1990）. “Oliver Heaviside”，The Engines of Our Ingenuity（《我们的创造引擎》），第426集，美国国家公共广播电台休斯敦KUHF-FM电台。
\item McGinty, Phil. “Oliver Heaviside”，刊于《Devon Life》，托贝图书馆服务。
\item Naughton, Russell. “Oliver W. Heaviside: 1850–1925”，刊于 Adventures in CyberSound（《赛博之声探险》）。
\item O’Connor, John J.; Robertson, Edmund F.，“Oliver Heaviside”，圣安德鲁斯大学数学史档案馆 MacTutor History of Mathematics Archive。
\item “Ron D.”（2007）.“Heaviside’s Operator Calculus”（《赫维赛德算子演算》）。
\item Eric W. Weisstein, “Heaviside, Oliver (1850–1925)”，刊于 Eric Weisstein’s World of Scientific Biography（《埃里克·魏斯坦科学传记大全》），Wolfram Media, Inc. 出版。
\end{itemize}
\subsubsection{档案收藏}
\begin{itemize}
\item Oliver Heaviside Selected Papers [microform], 1874–1922（《奥利弗·赫维赛德精选论文[缩微胶片]，1874–1922》），尼尔斯·玻尔图书馆与档案馆。
\end{itemize}